\subsubsection{Benchmark Bracket and Its Interpretation}
\label{benchmark-bracket}

\textit{Proposition (receiving-state bracket under independent primitives).}
Let $G$ be a finite DAG and let all initial-activation and edge-transmission
variables in Equation~(3-10) be mutually independent, with supplied parameters
in $[0,1]$. On the same graph and inputs, define recursive maximum and
clipped-additive states by
\[
M_j=\max\left(B_j,\max_{i\in P(j)}w_{ij}M_i\right),\qquad
H_j=\min\left(1,B_j+\sum_{i\in P(j)}w_{ij}H_i\right),
\eqtag{3-12}
\]
where the incoming maximum is zero for an empty predecessor set. Let $R_j$
be ordinary GBBRPM and $\pi_j=\Pr(X_j=1)$ under the benchmark. Let $T_j$
be the exact activation probability on the target-specific restricted
in-arborescence defined in Section~\ref{reference-comparators}. It will
retain a consistent maximum-weight path from every positive-path ancestor,
the original local inputs and weights, and target $j$ itself; no path-weight
threshold will be applied. Then
\[
M_j\le T_j\le\pi_j\le R_j\le H_j\qquad\text{for every }j\in V.
\eqtag{3-13}
\]

\textit{Proof.} In a DAG, $X_i$ cannot depend on the primitive $E_{ij}$,
and no incoming activation at $j$ can depend on $Z_j$. Independence of
the primitives therefore gives
$\Pr(X_i\land E_{ij})=w_{ij}\pi_i$. Since activation at $j$ includes
its own initial event and every successful incoming event,
$\pi_j\ge\max(B_j,\max_i w_{ij}\pi_i)$. A predecessor-free node has
$M_j=\pi_j=B_j$; induction over topological order establishes the lower bound.

The maximum recursion also has the path representation
$M_j=\max_{a,\gamma:a\rightarrow j}B_a\prod_{e\in\gamma}w_e$,
including the zero-length path at $j$. The restricted tree retains a
maximum-weight path for each positive-path ancestor by consistent next-edge
choices. Activation at $a$ through its own local event followed by successful
transmission along that retained path is sufficient to activate $j$ in the
tree. Its probability is $B_a\prod_{e\in\gamma}w_e$ by primitive independence.
The tree probability is at least each such path contribution and $B_j$, so
$M_j\le T_j$. Ancestors without a positive path contribute zero to $M_j$.
Couple the restricted and full processes using the same local and edge draws.
Deleting edges cannot create an activation, hence $T_j\le\pi_j$.
On the target-directed tree, incoming ancestor subtrees have disjoint
primitives; complement-product recursion therefore evaluates $T_j$ exactly.
This argument relies on the stated unthresholded, path-preserving construction;
arbitrary pruning need not retain the $T_j\ge M_j$ guarantee.

For the upper bound, let $F_{ij}$ denote failure of incoming activation
through $(i,j)$, that is, the complement of $\{X_i=1,E_{ij}=1\}$.
Every $X_i$ is an increasing Boolean function of the primitive variables,
so every $F_{ij}$ is decreasing. Positive association under a product
Bernoulli measure, the Harris/FKG principle, gives
$\Pr(\bigcap_iF_{ij})\ge\prod_i\Pr(F_{ij})$. This established principle
is also used for cascade upper bounds by \citet{abbe2017,abbe2020}; it is not a
new correlation inequality. Since $Z_j$ is independent of the incoming
events,
\[
\pi_j=1-(1-B_j)\Pr\left(\bigcap_{i\in P(j)}F_{ij}\right)
\le1-(1-B_j)\prod_{i\in P(j)}(1-w_{ij}\pi_i).
\eqtag{3-14}
\]
Complement product is increasing in each predecessor marginal.
Starting from $\pi_j=R_j=B_j$ at predecessor-free nodes, topological
induction yields $\pi_j\le R_j$.

Finally, for nonnegative inputs in $[0,1]$,
$1-\prod_k(1-x_k)\le\min(1,\sum_kx_k)$. Applying this inequality to
the local term and incoming contributions, followed by topological
induction and monotonicity of clipped addition, gives $R_j\le H_j$.
Primitive probabilities of zero or one follow directly or by continuity.
This completes the argument.

The proposition will be presented as a specialization of established
association and bounded-aggregation arguments, not as a verified first
theorem about cascade approximation. Abbe et al.'s journal treatment
also bounds individual activation probabilities before summing them into
seed-set influence. Its upper method uses length-indexed messages; its
DAG-based lower method uses a different incoming-message recurrence
\citep{abbe2020}. Section~\ref{cascade-bounds-prior-art} compares the
constructions and limits. Independent local activation
can also be represented by one certainly active auxiliary source with
independent edges of probabilities $B_j$ to all original nodes. That
construction facilitates comparison with fixed-seed literature; it will
not conceal the model's mathematical precedents.

For this benchmark, score discrepancy is one-sided: $R_j-\pi_j\ge0$.
Its scalar certificate is $0\le R_j-\pi_j\le R_j-M_j$.
When the tree comparator is computed, the tighter certificate is
$0\le R_j-\pi_j\le R_j-T_j\le R_j-M_j$.
Computing $T_j$ requires tree construction and evaluation for each requested
target; the tighter certificate is not free relative to the scalar pass.
Its reporting will reuse the already computed comparator result.
GBBRPM is no looser than recursive clipped addition as a nodewise upper
bound; this ordering need not be rediscovered experimentally. Experiments
will measure bound width, actual slack, conditions of equality, and
ranking or selection effects. A smaller pointwise score error does not
guarantee better rank order. The condition $T_a>R_b$ certifies
$\pi_a>\pi_b$ when tree bounds are available; $M_a>R_b$ remains a cheaper
sufficient certificate. Overlapping brackets remain inconclusive,
not evidence of a reversed ordering.

These guarantees concern only the stated independent-primitive benchmark.
They do not certify probabilities or physical bounds for the domain index,
dependent initial events, correlated transmission primitives, feedback, or
other cascade mechanisms. The ancestry-aware diagnostic will require its
own analysis and will not inherit the GBBRPM upper-bound guarantee.
